\documentclass[11pt,a4paper]{article}

\usepackage[margin=2.5cm]{geometry}
\usepackage{amsmath}
\usepackage{amssymb}
\usepackage{graphicx}
\usepackage{booktabs}
\usepackage[colorlinks=true,linkcolor=blue,citecolor=blue,urlcolor=blue]{hyperref}

\newcommand{\dq}{\ensuremath{dq}}

\title{Even Harmonics in Three-Phase Active Front-End Rectifiers:\\
Origins, Mechanisms, and Mitigation}
\author{Davide Bagnara}
\date{\today}

\begin{document}

\maketitle

\begin{abstract}
Even-order current harmonics are non-characteristic for a three-phase active
front-end (AFE) rectifier: under ideal half-wave symmetry they cannot exist.
Their appearance in the field is therefore always a symptom of a symmetry-breaking
mechanism in the measurement chain, in the power stage, in the modulator, or in
the grid itself. This note reviews the physical and control-theoretic origins of
even harmonics in grid-connected voltage-source converters, with emphasis on the
DC-offset\,$\rightarrow$\,second-harmonic chain, asymmetric switching delays,
DC injection with transformer half-cycle saturation, asynchronous low-ratio
modulation, and interactions with an unbalanced or pre-distorted grid. The
relevant emission and compatibility limits from power-system standards are
summarized, and practical mitigation measures are given.
\end{abstract}

\section{Introduction}

A three-phase AFE rectifier operated with a symmetric modulation strategy and an
ideal power stage draws line currents containing only the characteristic
harmonic families of a balanced six-switch converter: the low-frequency
characteristic orders $h = 6k \pm 1$ ($k = 1, 2, \dots$) when the converter is
current-source-like or line-commutated, and, for a PWM voltage-source converter,
sidebands clustered around integer multiples of the switching frequency
\cite{mohan,holmes}. Even-order components ($h = 2, 4, 6, \dots$) and a DC
component are \emph{non-characteristic}: they are forbidden by the waveform
symmetry of the ideal converter.

In practice, small even-order components---most often a second harmonic
accompanied by a DC component---are routinely measured on high-power AFE
installations. Because they are non-characteristic, they carry diagnostic
information: each possible origin corresponds to a specific defect or asymmetry.
This note collects these origins in one place and links the converter-level
(power electronics) view to the network-level (power systems) view, where even
harmonics interact with transformer saturation, protection relaying, and
compatibility limits.

\section{Half-Wave Symmetry and Non-Characteristic Harmonics}

\subsection{Fourier condition}

A periodic waveform $f(t)$ with period $T$ has \emph{half-wave symmetry} if
%
\begin{equation}
f\!\left(t + \tfrac{T}{2}\right) = -f(t).
\label{eq:hws}
\end{equation}
%
Writing the Fourier series $f(t) = \sum_n \left[a_n \cos n\omega_1 t
+ b_n \sin n\omega_1 t\right]$ and imposing \eqref{eq:hws} gives
%
\begin{equation}
a_n = b_n = 0 \qquad \text{for all even } n \text{ (including } n=0\text{)}.
\end{equation}
%
Hence a waveform satisfying \eqref{eq:hws} contains \emph{no} DC component and
\emph{no} even harmonics. In a balanced three-phase converter the modulating
signals, the switching pattern, and the resulting line-to-line voltages all
satisfy \eqref{eq:hws} in the ideal case; the line currents inherit the
property through the linear grid impedance. Every observed even harmonic is
therefore the signature of a violated symmetry: the positive and negative
half-cycles are no longer mirror images of each other. DC and even orders are
two faces of the same defect and almost always appear together
\cite{arrillaga,wakileh}.

\subsection{Frequency mapping between \texorpdfstring{$abc$}{abc} and
\texorpdfstring{\dq}{dq} frames}

Many of the mechanisms below are best understood through the frequency mapping
of the synchronous ($\dq$) reference frame rotating at the fundamental
angular frequency $\omega_1$. A positive-sequence component at frequency
$\omega$ in the stationary frame appears at $\omega - \omega_1$ in the
$\dq$ frame; a negative-sequence component at $\omega$ appears at
$\omega + \omega_1$. The mappings relevant to this note are:
%
\begin{center}
\begin{tabular}{lll}
\toprule
$abc$ (stationary) & sequence & $\dq$ (synchronous) \\
\midrule
DC offset                & --            & $1\omega_1$ ripple \\
fundamental              & positive      & DC \\
fundamental              & negative      & $2\omega_1$ ripple \\
2nd harmonic             & negative      & $3\omega_1$ ripple \\
2nd harmonic             & positive      & $1\omega_1$ ripple \\
5th (neg.) and 7th (pos.)& --            & $6\omega_1$ ripple \\
\bottomrule
\end{tabular}
\end{center}
%
Two entries deserve attention. First, a DC offset in the measured $abc$
currents and a positive-sequence second harmonic \emph{both} map to a
$1\omega_1$ ripple in $\dq$; a synchronous-frame current controller cannot
distinguish them and, worse, converts one into the other, as shown in
Section~\ref{sec:offset}. Second, grid unbalance (negative-sequence
fundamental) maps to $2\omega_1$ in $\dq$, which explains the well-known
$2\omega_1$ DC-bus ripple under unbalanced conditions \cite{yazdani,teodorescu}.

\section{Origin Mechanisms}

\subsection{DC offset in the current (or voltage) measurements}
\label{sec:offset}

The most common single cause of second-harmonic emission in AFE converters is a
DC offset in the current-sensing chain (Hall sensor, burden resistor, analog
conditioning, or ADC). The mechanism is a closed-loop chain:
%
\begin{enumerate}
\item A sensor offset adds a spurious DC component to the measured $abc$
currents.
\item After the Park transformation, this DC error appears as a
$1\omega_1$ ripple superimposed on the $\dq$ current feedback.
\item The current controller, whose bandwidth normally exceeds $\omega_1$ by a
comfortable margin, reacts to this fictitious ripple and injects a real
$1\omega_1$ component into the $\dq$ voltage reference.
\item Transformed back to $abc$, a $1\omega_1$ component in $\dq$ produces a
DC component \emph{and} a second-harmonic component in the actual converter
voltages and hence in the line currents.
\end{enumerate}
%
The chain \emph{offset $\rightarrow$ DC $+$ 2nd harmonic} is so repeatable that
the simultaneous presence of a DC component and a second harmonic in the line
current is practically a diagnostic fingerprint of measurement offset. The
analogous mechanism in motor drives, where offset and gain errors of the current
sensors produce first- and second-harmonic torque ripple, was analyzed by Chung
and Sul \cite{chungsul}; the AFE case is formally identical with the grid
replacing the machine. Gain \emph{mismatch} between phase sensors, by contrast,
creates a spurious negative-sequence fundamental ($2\omega_1$ in $\dq$) and
leads to odd, not even, distortion---the distinction between offset errors
(even orders) and gain errors (odd orders) is useful in commissioning.

Voltage-measurement offsets act similarly through the PLL and the feedforward
path: an offset on the grid-voltage measurement modulates the PLL angle at
$1\omega_1$, which again produces DC and second-harmonic components in the
synthesized converter voltage.

\subsection{Asymmetric switching delays and dead time}

Every switching event of a converter leg introduces a volt-second error with
respect to the ideal waveform. With dead time $t_d$, turn-on delay
$t_{on}$, and turn-off delay $t_{off}$, the average voltage error of one leg
over a switching period $T_{sw}$ has magnitude
%
\begin{equation}
\Delta V \;=\; \frac{t_d + t_{on} - t_{off}}{T_{sw}}\, V_{dc}
\;\cdot\; \operatorname{sign}(i),
\label{eq:deadtime}
\end{equation}
%
where the sign follows the leg current polarity. As long as the error is
\emph{identical} for both current polarities---i.e., the top and bottom devices
and their gate drivers are matched---the distortion in \eqref{eq:deadtime} is an
odd function of the current and preserves half-wave symmetry: the classical
dead-time distortion therefore contains only odd harmonics (predominantly 5th
and 7th in a three-phase system, since triplen components cancel)
\cite{jeong,choi}.

Even harmonics appear when the error is \emph{different} between the two
half-cycles:
%
\begin{itemize}
\item unequal propagation delays of the top and bottom gate-driver channels;
\item asymmetric dead-time insertion (different $t_d$ for rising and falling
edges);
\item parameter spread between the top and bottom devices ($t_{on}$,
$t_{off}$, $V_{CE,\mathrm{sat}}$ or $R_{DS,\mathrm{on}}$);
\item unequal snubber or gate-resistor values after service or rework.
\end{itemize}
%
Any such top/bottom asymmetry adds a polarity-independent term to
\eqref{eq:deadtime}, i.e.\ a DC voltage error per leg, plus even-order
distortion. At high power levels with low switching frequency the relative
weight of a fixed nanosecond-scale mismatch grows with $1/T_{sw}$ being small,
but the volt-second error per event is applied fewer times per fundamental
period; in practice, mismatches of tens to hundreds of nanoseconds between
driver channels are measurable in the line-current spectrum of megawatt
converters.

\subsection{DC injection and transformer half-cycle saturation}

Whatever its origin, a DC component injected by the converter flows through the
line-side magnetics. A transformer offers essentially no impedance to DC other
than its winding resistance, so even a small DC voltage bias drives a
substantial DC magnetizing current. The flux is displaced from its symmetric
trajectory and the core saturates on one half-cycle only. The resulting
magnetizing current is a train of unipolar peaks, extremely rich in
second-harmonic content---the same phenomenon as transformer inrush, which is
precisely why differential protection relays use second-harmonic restraint to
distinguish inrush from internal faults \cite{guzman}.

Two aspects make this mechanism particularly relevant:
%
\begin{itemize}
\item \textbf{Self-reinforcement.} Half-cycle saturation itself produces even
harmonics and DC in the current, which can interact with the converter control
and with other transformers in the network, sustaining or amplifying the
asymmetry.
\item \textbf{Regulatory limits on DC injection.} Interconnection standards
constrain the DC current a converter may inject: IEEE Std 1547 limits it to
0.5\,\% of rated current \cite{ieee1547}. The limit exists largely because of
the saturation mechanism above, including its effect on distribution
transformers and on residual-current devices.
\end{itemize}
%
The same physics operates at network scale during geomagnetic disturbances:
geomagnetically induced currents (GIC) are quasi-DC currents that drive
transmission transformers into half-cycle saturation, generating large even
(and odd) harmonic currents and reactive-power consumption across entire
networks \cite{pirjola}. An AFE installed on a grid affected by transformer
saturation---whatever the cause---will see a background second-harmonic voltage
and respond to it (Section~\ref{sec:background}).

\subsection{Low and asynchronous modulation ratio}

At the power levels typical of megawatt AFE converters the switching frequency
is low and the frequency-modulation ratio $m_f = f_{sw}/f_1$ may be below
$\sim 21$. In this regime the classical PWM harmonic theory \cite{holmes}
prescribes \emph{synchronous} modulation with $m_f$ an odd integer (ideally an
odd multiple of three): the carrier is phase-locked to the fundamental, the
pulse pattern repeats identically every fundamental period, and half-wave
symmetry of the switched voltage is guaranteed, so no even harmonics and no
sub- or interharmonics are produced.

If the modulation is asynchronous (fixed carrier not locked to the grid) with a
small non-integer $m_f$, the pulse pattern drifts with respect to the
fundamental. The switched waveform is no longer periodic over one fundamental
cycle, half-wave symmetry is lost on a cycle-by-cycle basis, and the spectrum
contains even harmonics, subharmonics, and interharmonics whose amplitude grows
as $m_f$ decreases. The remedy is structural, not a tuning matter: lock the
carrier to the PLL angle and enforce an odd (preferably triplen-odd) $m_f$, or
adopt an optimized synchronous pattern (selective harmonic elimination) that
embeds half-wave and quarter-wave symmetry by construction.

\subsection{Grid unbalance and control interaction}

Grid unbalance is frequently blamed for even harmonics, but the mapping must be
kept precise. A negative-sequence fundamental component (unbalance) appears as
a $2\omega_1$ ripple in the $\dq$ frame and as a $2\omega_1$ power/voltage
ripple on the DC bus. If the DC-bus voltage controller is fast enough to chase
this ripple, it modulates the $d$-axis current reference at $2\omega_1$; back
in the stationary frame this produces a \emph{third} harmonic (positive
sequence) and an additional negative-sequence fundamental---odd components, not
even ones \cite{yazdani,teodorescu}.

The correct attribution is therefore:
%
\begin{itemize}
\item \textbf{Offset-type asymmetries} (sensor offsets, top/bottom leg
asymmetry, DC injection) $\rightarrow$ DC and \emph{even} harmonics;
\item \textbf{Gain-type asymmetries and unbalance} (sensor gain mismatch,
negative-sequence grid voltage, fast DC-bus loop) $\rightarrow$
negative-sequence fundamental and \emph{odd} harmonics (3rd in particular).
\end{itemize}
%
Keeping the two families separate greatly accelerates root-cause analysis on
site: the parity of the dominant non-characteristic harmonic identifies the
class of defect before any hardware is opened.

\subsection{Background distortion of the grid}
\label{sec:background}

Finally, the converter may simply be responding to a grid that already carries
even-harmonic voltage distortion. Sources on the network side include
half-wave-rectifier loads, cycloconverters and other asymmetrically fired
line-commutated converters, transformer energization (inrush) events nearby,
arc furnaces during unstable arcing, and GIC-driven transformer saturation
\cite{arrillaga,bollen,taskforce}. An AFE presents a low impedance around the
fundamental and its controller bandwidth; a background second-harmonic voltage
therefore drives a second-harmonic current through the converter admittance
unless the control actively rejects it. In emission assessments this
distinction matters: the measured even-harmonic current at the point of common
coupling is the superposition of the converter's own emission and its response
to background voltage distortion, and the two must be separated before
compliance is judged \cite{iec61000_3_6}.

\section{Standards Perspective}

Power-system standards treat even harmonics more strictly than odd ones,
precisely because they are non-characteristic and symptomatic of defects:
%
\begin{itemize}
\item \textbf{IEEE Std 519} limits even current harmonics to 25\,\% of the
corresponding odd-harmonic limits in each frequency range, and additionally
restricts DC injection \cite{ieee519}.
\item \textbf{IEC 61000-3-6} provides planning levels and the methodology for
apportioning emission limits to distorting installations connected to MV, HV,
and EHV systems, with dedicated (lower) indicative planning levels for even
orders \cite{iec61000_3_6}.
\item \textbf{EN 50160} specifies the voltage-quality characteristics at the
supply terminals of public networks, including individual even-harmonic voltage
limits that are markedly lower than those of neighboring odd orders (e.g.\
2\,\% for $h=2$ and 1\,\% for $h=4$, against 5--6\,\% for the low odd orders)
\cite{en50160}.
\item \textbf{IEEE Std 1547} limits DC injection of distributed resources to
0.5\,\% of rated output current, addressing the transformer-saturation chain
described above \cite{ieee1547}.
\end{itemize}

\section{Mitigation}

The countermeasures follow directly from the mechanisms:
%
\begin{itemize}
\item \textbf{Offset auto-calibration.} Measure and store the ADC/sensor
offsets at every start-up with the converter gated off, and track slow offset
drift online (e.g.\ by averaging the measured currents over an integer number
of fundamental periods during operation).
\item \textbf{Resonant compensation.} Add a resonant term at $1\omega_1$ in the
$\dq$-frame current controller (equivalently, at DC and $2\omega_1$ in the
stationary frame). This simultaneously rejects residual measurement offset
effects, DC injection, and background second-harmonic voltage. Repetitive
control achieves the same result for all harmonic orders at once.
\item \textbf{Gate-path symmetry.} Match propagation delays of top and bottom
driver channels, use symmetric dead-time insertion, and apply dead-time /
delay compensation based on current polarity. Verify after any driver or
device replacement.
\item \textbf{Synchronous modulation.} At low $m_f$, lock the carrier to the
grid angle and use an odd (triplen) integer ratio, or an offline-optimized
symmetric pattern.
\item \textbf{DC suppression loop.} For transformer-coupled converters, close a
dedicated slow loop on the estimated DC component of the line current (or on
the transformer magnetizing asymmetry) and trim the modulator offset
accordingly \cite{buticchi}.
\item \textbf{Decoupled sequence control.} Under unbalanced grid conditions,
use dual (positive/negative) sequence current control and limit the DC-bus
voltage loop bandwidth well below $2\omega_1$, so that unbalance is not
converted into additional low-order distortion.
\end{itemize}

\section{Conclusion}

Even harmonics in a three-phase AFE are never intrinsic: they always indicate a
broken half-wave symmetry. The practically dominant chain is measurement DC
offset producing DC plus second harmonic through the synchronous-frame current
controller; the remaining mechanisms---top/bottom leg asymmetry, DC injection
with transformer half-cycle saturation, asynchronous low-ratio modulation, and
background network distortion---each leave a distinguishable signature. Because
offset-type defects generate even orders while gain-type defects and unbalance
generate odd orders, the parity of the observed non-characteristic harmonics is
itself a powerful diagnostic. Standards on both the emission side (IEEE 519,
IEC 61000-3-6) and the compatibility side (EN 50160) reflect this by assigning
even orders limits far below those of odd orders, and by constraining DC
injection directly.

\begin{thebibliography}{99}

\bibitem{mohan}
N.~Mohan, T.~M. Undeland, and W.~P. Robbins,
\emph{Power Electronics: Converters, Applications, and Design}, 3rd~ed.
Hoboken, NJ: Wiley, 2003.

\bibitem{holmes}
D.~G. Holmes and T.~A. Lipo,
\emph{Pulse Width Modulation for Power Converters: Principles and Practice}.
Piscataway, NJ: IEEE Press / Wiley-Interscience, 2003.

\bibitem{yazdani}
A.~Yazdani and R.~Iravani,
\emph{Voltage-Sourced Converters in Power Systems: Modeling, Control, and
Applications}. Hoboken, NJ: Wiley / IEEE Press, 2010.

\bibitem{teodorescu}
R.~Teodorescu, M.~Liserre, and P.~Rodr\'iguez,
\emph{Grid Converters for Photovoltaic and Wind Power Systems}.
Chichester, UK: Wiley / IEEE Press, 2011.

\bibitem{chungsul}
D.-W. Chung and S.-K. Sul,
``Analysis and compensation of current measurement error in vector-controlled
AC motor drives,'' \emph{IEEE Trans. Ind. Appl.}, vol.~34, no.~2,
pp. 340--345, Mar./Apr. 1998.

\bibitem{jeong}
S.-G. Jeong and M.-H. Park,
``The analysis and compensation of dead-time effects in PWM inverters,''
\emph{IEEE Trans. Ind. Electron.}, vol.~38, no.~2, pp. 108--114, Apr. 1991.

\bibitem{choi}
J.-W. Choi and S.-K. Sul,
``Inverter output voltage synthesis using novel dead time compensation,''
\emph{IEEE Trans. Power Electron.}, vol.~11, no.~2, pp. 221--227, Mar. 1996.

\bibitem{buticchi}
G.~Buticchi, E.~Lorenzani, and G.~Franceschini,
``A DC offset current compensation strategy in transformerless grid-connected
power converters,'' \emph{IEEE Trans. Power Del.}, vol.~26, no.~4,
pp. 2743--2751, Oct. 2011.

\bibitem{arrillaga}
J.~Arrillaga and N.~R. Watson,
\emph{Power System Harmonics}, 2nd~ed. Chichester, UK: Wiley, 2003.

\bibitem{wakileh}
G.~J. Wakileh,
\emph{Power Systems Harmonics: Fundamentals, Analysis and Filter Design}.
Berlin, Germany: Springer, 2001.

\bibitem{bollen}
M.~H.~J. Bollen,
\emph{Understanding Power Quality Problems: Voltage Sags and Interruptions}.
Piscataway, NJ: IEEE Press, 2000.

\bibitem{taskforce}
Task Force on Harmonics Modeling and Simulation,
``Modeling and simulation of the propagation of harmonics in electric power
networks---Part I: Concepts, models, and simulation techniques,''
\emph{IEEE Trans. Power Del.}, vol.~11, no.~1, pp. 452--465, Jan. 1996.

\bibitem{guzman}
A.~Guzm\'an, S.~Zocholl, G.~Benmouyal, and H.~J. Altuve,
``A current-based solution for transformer differential protection---Part I:
Problem statement,'' \emph{IEEE Trans. Power Del.}, vol.~16, no.~4,
pp. 485--491, Oct. 2001.

\bibitem{pirjola}
R.~Pirjola,
``Geomagnetically induced currents during magnetic storms,''
\emph{IEEE Trans. Plasma Sci.}, vol.~28, no.~6, pp. 1867--1873, Dec. 2000.

\bibitem{ieee519}
\emph{IEEE Standard for Harmonic Control in Electric Power Systems},
IEEE Std 519-2022, 2022.

\bibitem{iec61000_3_6}
\emph{Electromagnetic Compatibility (EMC)---Part 3-6: Limits---Assessment of
Emission Limits for the Connection of Distorting Installations to MV, HV and
EHV Power Systems}, IEC/TR 61000-3-6, 2008.

\bibitem{en50160}
\emph{Voltage Characteristics of Electricity Supplied by Public Electricity
Networks}, EN 50160, CENELEC, 2010.

\bibitem{ieee1547}
\emph{IEEE Standard for Interconnection and Interoperability of Distributed
Energy Resources with Associated Electric Power Systems Interfaces},
IEEE Std 1547-2018, 2018.

\end{thebibliography}

\end{document}
